%% Proposed addition to Section 6 of the manuscript, to follow Section 6.5 (Balzer's Conditions).
%% It uses the macros and theorem environments of paper.tex and refers to these labels of the
%% manuscript: sec:balzer, prop:balzer, lem:lightcone, lem:agree, lem:return, lem:pumping, thm:left.
%% The manuscript itself is not changed by this file.

\subsection{Stable Generals}\label{sec:stable}

\Cref{prop:balzer} uses Balzer's four conditions together. Far less is needed. The
first condition alone excludes five states, and so does its weak reading
together with the image condition for the neighbourhoods without a border. The
second statement contains the form in which Yun\`es states the conditions, which
\cref{sec:balzer} left open.

\begin{theorem}[computer-assisted]\label{thm:stable}
No five-state rule that satisfies \textup{(B1)} synchronizes all lines of
lengths $2,\dots,16$ in minimal time. The lengths up to $14$ do not suffice to
prove this.
\end{theorem}

\begin{proposition}[computer-assisted]\label{prop:weakimage}
For an involution $I$ of $W=\{\Lq,\Gq,\Aq,\Bq\}$, extended by $I(\Fq)=\Fq$, let
$N^\ast_I$ be the largest $N$ such that some five-state rule that satisfies
\textup{(B1)} in the weak reading and
\begin{equation}\label{eq:B2o}
\delta(I(z),I(y),I(x))=I(\delta(x,y,z))\qquad\text{for all }x,y,z\in W
\end{equation}
synchronizes all lines of lengths $2,\dots,N$ in minimal time. Then
$N^\ast_I\in\{12,13\}$ for the identity, $N^\ast_I\in\{10,11\}$ for
$(\Aq\,\Bq)$, $(\Lq\,\Aq)$ and $(\Lq\,\Bq)$, $N^\ast_I=6$ for $(\Gq\,\Aq)$,
$(\Gq\,\Bq)$, $(\Lq\,\Aq)(\Gq\,\Bq)$ and $(\Lq\,\Bq)(\Gq\,\Aq)$, and $N^\ast_I=2$
for $(\Lq\,\Gq)$ and $(\Lq\,\Gq)(\Aq\,\Bq)$. In particular, no five-state rule that
satisfies these two conditions for some involution synchronizes all lines of
lengths $2,\dots,14$ in minimal time.
\end{proposition}

\Cref{prop:weakimage} assumes nothing about the neighbourhoods with image
$\Fq$, about the values $\delta(\Gq,v,\Gq)$, or about the neighbourhoods with a
border, and $I$ may move $\Gq$. With $I$ the identity, its hypotheses hold for the rule
\texttt{balzer\_noborder12} of \cref{sec:balzer}, which gives the lower bound
$12$. For the class of \cref{thm:stable}, the largest $N$ such that some rule
synchronizes the lines of lengths $2,\dots,N$ is $14$ or $15$. The search below
found a rule, \texttt{stable\_generals\_2-14},
that satisfies (B1), gives $\Fq$ exactly on $(\Gq,\Gq,\Gq)$, $(\ast,\Gq,\Gq)$
and $(\Gq,\Gq,\ast)$, and synchronizes the lines of lengths $2,\dots,14$ and no
other line of length at most $120$ (direct simulation). On its half-line the
cells $2,\dots,6$ have period $2$ in time from time $6$ on and every cell $i\ge7$
is constant from time $i-1$ on (simulation for $t<150$), which is the kind of
half-line that \cref{thm:left} excludes for a solution. Up to the transitions
that are used neither in the lines of lengths $2,\dots,14$ nor in the half-line
up to the anti-diagonal $27$, and up to the exchange of $\Aq$ and $\Bq$, it is one
of exactly three models of the formula $F_{14}$ of the proof below: the other
two differ from it in $\delta(\Aq,\Aq,\Aq)$, $\delta(\Gq,\Aq,\Aq)$ and in a few
transitions that it does not use, and the refutation of $F_{14}$ with the
node $\gamma^\ast$ and three clauses that exclude these models has a proof
verified by both checkers.

\paragraph{Consequences of a stable general}
A \emph{barrier} of the line of length $n$ at time $t$ is a border (position
$0$ or $n+1$) or a cell in state $\Gq$ at time $t$.

\begin{lemma}\label{lem:fire}
Let $\delta$ satisfy \textup{(B1)}.
\begin{enumerate}[(1)]
\item If $\delta$ synchronizes the lines of lengths $2$ and $3$ at times $2$ and
$4$, then $\delta(\Gq,\Lq,\ast)=\Gq$ and
$\delta(\ast,\Gq,\Gq)=\delta(\Gq,\Gq,\ast)=\delta(\Gq,\Gq,\Gq)=\Fq$.
\item If moreover $\delta$ synchronizes the line of length $n\ge2$ at time
$2n-2$, then
\begin{enumerate}[(a)]
\item $C_n(t,1)=\Gq$ for all $t\le2n-3$;
\item $C_n(2n-3,i)=\Gq$ for $1\le i\le n$;
\item at no time $t\le2n-4$ are three adjacent positions barriers;
\item if $C_n(t,i)=\Gq$ and $t\le2n-4$, then $C_n(t+1,i)=\Gq$.
\end{enumerate}
\end{enumerate}
\end{lemma}

\begin{proof}
(1) In the line of length $2$, $C_2(1,1)=\delta(\ast,\Gq,\Lq)=\Gq$ by (B1), and
$\Fq=C_2(2,1)=\delta(\ast,\Gq,x)$ with $x:=C_2(1,2)=\delta(\Gq,\Lq,\ast)\in W$. By
(B1) this value is $\Gq$ unless $x=\Gq$. Hence $x=\Gq$ and
$\delta(\ast,\Gq,\Gq)=\Fq$, and then $\Fq=C_2(2,2)=\delta(\Gq,\Gq,\ast)$. In the
line of length $3$ let $s$ be the first time at which the cell $2$ is in state
$\Gq$; it exists, because otherwise the cell $1$ would be in state $\Gq$ for
ever, by (B1). The cell $1$ is in state $\Gq$ up to time $s$ and in state
$\delta(\ast,\Gq,\Gq)=\Fq$ at time $s+1$, so $s+1=4$. Then
$\Fq=C_3(4,2)=\delta(\Gq,\Gq,z)$ with $z:=C_3(3,3)\in W$, and (B1) gives $z=\Gq$
and $\delta(\Gq,\Gq,\Gq)=\Fq$.

(2c) If three adjacent positions are barriers at time $t$, the middle one is a
cell in state $\Gq$, because for $n\ge2$ the two borders are not adjacent to a
common cell; by (1) this cell is in state $\Fq$ at time $t+1$, so
$t+1\ge2n-2$. (2d) By (2c) the two neighbours of the cell are not both
barriers, and (B1) applies. (2a) By (2d) and induction on $t$, starting from
$C_n(0,1)=\Gq$. (2b) By (2a), $\Fq=C_n(2n-2,1)=\delta(\ast,\Gq,y)$ with
$y:=C_n(2n-3,2)\in W$, so $y=\Gq$ by (B1). If $C_n(2n-3,i)=\Gq$ and $i<n$, then
$\Fq=C_n(2n-2,i)=\delta(x,\Gq,z)$ with $z:=C_n(2n-3,i+1)\in W$, and (B1) gives
$z=\Gq$.
\end{proof}

So a rule with (B1) produces $\Fq$, on every line that it synchronizes, only at
time $2n-2$, from the row $\Gq^n$ and the three neighbourhoods of (B3). The
strong reading of (B3) can therefore be enforced without changing the line of
any synchronized length, by redefining $\delta(x,y,z):=y$ for the other
neighbourhoods with image $\Fq$; these have $y\ne\Gq$ by (B1), so (B1) is
preserved.

\begin{lemma}[walls]\label{lem:wall}
Let $N\ge3$, and let $\delta$ satisfy \textup{(B1)} and synchronize the lines of
lengths $2,\dots,N$ in minimal time. Then $C_\infty(t,1)=\Gq$ for $t\le2N-3$, and
$C_\infty(t,i)\ne\Gq$ whenever $i\ge2$ and $t+i\le2N-4$.
\end{lemma}

\begin{proof}
The first claim is \cref{lem:fire}\,(2a) for $n=N$ and \cref{lem:agree}. Let
$p\ge2$ and $t_0+p\le2N-4$, and suppose that $C_\infty(t_0,p)=\Gq$. Put
$n:=N-1$, so that $t_0+p\le2n-2$ and $t_0\le2n-4$. By \cref{lem:agree},
$C_n(t_0,i)=C_{n+1}(t_0,i)=C_\infty(t_0,i)$ for $i\le p$. By
\cref{lem:fire}\,(2d), applied to both lengths, the cell $p$ is in state $\Gq$ at
the times $t_0,\dots,2n-3$ in both lines. The states of the cells
$1,\dots,p-1$ at time $t+1$ are determined by the states of the cells
$1,\dots,p$ at time $t$; by induction on $t$, the cells $1,\dots,p-1$ have the
same states in both lines at all times $t_0\le t\le2n-2$. Hence
$C_{n+1}(2n-2,1)=C_n(2n-2,1)=\Fq$, although $2n-2<2(n+1)-2$.
\end{proof}

Thus a cell in state $\Gq$ is a wall that no information crosses, a wall other
than the first cell can appear only after the echo of the right end has passed,
and the half-line of a solution with (B1) would be a diagram over
$W':=W\setminus\{\Gq\}$ next to a constant first cell. The cells between two
walls evolve on their own, and they cannot do so for long.

\begin{lemma}[gaps]\label{lem:gap}
Let $\delta$ satisfy \textup{(B1)} and synchronize the lines of lengths $2$, $3$
and $n$ in minimal time. Suppose that at a time $t\le2n-4$ two barriers of the
line of length $n$ are at distance $\ell+1\ge2$ and no position between them is
a barrier. Then
\[
2n-3-t\le M(\ell),\qquad M(1):=|W'|,\quad M(\ell):=|W'|^{\ell}+M(\ell-1)\ \
(\ell\ge2).
\]
Consequently, if $2n-3-t>M(\ell)$, no two barriers of the line are at a
distance $d$ with $2\le d\le\ell+1$ at time $t$.
\end{lemma}

\begin{proof}
By \cref{lem:fire}\,(2d) the two barriers stay barriers up to time $2n-3$, so at
the times $t,\dots,2n-3$ the $\ell$ cells between them are driven by themselves
and two constant neighbours. As long as none of these cells is in state $\Gq$,
their configuration lies in $W'^{\ell}$ and does not repeat: after a repetition
it would be periodic up to time $2n-3$, against \cref{lem:fire}\,(2b). Hence
there is $1\le m\le|W'|^{\ell}$ such that one of the $\ell$ cells is in state
$\Gq$ at time $t+m\le2n-3$. If $t+m=2n-3$, then
$2n-3-t\le|W'|^{\ell}\le M(\ell)$. Otherwise $t+m\le2n-4$; by
\cref{lem:fire}\,(2c) not all of the $\ell$ cells are in state $\Gq$ at time
$t+m$, which is impossible for $\ell=1$, and for $\ell\ge2$ some two barriers at
time $t+m$ are at a distance $\ell'+1$ with $1\le\ell'\le\ell-1$ and have no
barrier between them. By induction, $2n-3-(t+m)\le M(\ell')\le M(\ell-1)$. For
the last statement, let two barriers be at a distance $d$ with
$2\le d\le\ell+1$ at time $t$. Not all positions between them are barriers, by
\cref{lem:fire}\,(2c); so two barriers at a distance $\ell'+1\le d$ have no
barrier between them, and $2n-3-t\le M(\ell')\le M(\ell)$.
\end{proof}

For five states, $M(1),M(2),M(3)=3,12,39$.

\paragraph{A case split over the half-line}
As single formulas, the classes of \cref{thm:stable,prop:weakimage} are hard
for the solvers that we used; they become easy once a beginning of the
half-line is fixed. Fix $N\ge3$ and put $M:=2N-1$. The formula $F^{(N)}$ of a
class of five-state rules has one-hot variables $[\delta(e)=d]$ for every
neighbourhood $e$ and state $d$, a one-hot vector of variables over $W$ for
every cell $(t,i)$ of the half-line with $t\ge1$, $1\le i\le t+1$ and
$t+i\le M$, and one for every cell $(t,i)\in R_n$ with $t\le2n-3$, for
$2\le n\le N$; the other cells of the lines are identified with cells of the
half-line (\cref{lem:agree}) or are constants. Its clauses say that
\begin{enumerate}[(F1)]
\item every cell with variables is the image under $\delta$ of the three cells
above it;
\item for $2\le n\le N$, every cell of the line of length $n$ at time $2n-3$ has
a neighbourhood with image $\Fq$;
\item $C_n(2n-1-i,i)\ne C_\infty(2n-1-i,i)$ for $2\le n\le N$ and $2\le i\le n$;
\item the front cells $(t,t+1)$, $1\le t\le N-1$, are not $\Lq$;
\item the conditions of the class hold.
\end{enumerate}

\begin{lemma}\label{lem:Fsound}
If a rule of the class synchronizes all lines of lengths $2,\dots,N+2$ in minimal
time, then its transitions and the cells of its lines and half-line satisfy
$F^{(N')}$ for every $3\le N'\le N$. If $N\ge5$, the same holds for a five-state
rule of the class that synchronizes all lines of lengths $2,\dots,N+1$.
\end{lemma}

\begin{proof}
It suffices to treat $N'=N$. The half-line cells with $t+i\le2N-1$ agree with
the line of length $N+1$, at times $t\le2N-2<2(N+1)-2$; so they are working
states, and (F1) holds for them and for the cells of the lines $n\le N$ before
time $2n-2$. (F2) is the synchronization of the lines $n\le N$. (F3) is
\cref{lem:return} for $n\le N$, which uses the line of length $n+1\le N+1$. (F4):
by \cref{lem:return} the front cell $(n-2,n-1)$ is not $\Lq$ if the lines of
lengths $n$ and $n+1$ are synchronized; with $n\le N+1$ this gives the cells
$(t,t+1)$ for $t\le N-1$. For the second statement only (F4) for $t=N-1$ has to
be shown again. The front cells $g_t:=C_\infty(t,t+1)$ satisfy $g_0=\Gq$ and
$g_{t+1}=\delta(g_t,\Lq,\Lq)$, by \cref{lem:lightcone}. The lines of lengths up to
$N+1$ give $g_t\ne\Lq$ for $t\le N-2$, in particular for $t\le3$. So
$g_0,\dots,g_3$ lie in the three-element set $W\setminus\{\Lq\}$, two of them are
equal, and the sequence is periodic from there on with values in
$W\setminus\{\Lq\}$.
\end{proof}

Order the half-line cells with variables by $(t,i)\prec(t',i')$ if $t+i<t'+i'$,
or $t+i=t'+i'$ and $t<t'$. The three cells above a cell precede it in this
order or are constants. A \emph{node} is a map $\gamma$ from a finite set of
neighbourhoods to $W$. It determines the states of an initial segment of the
cells; let $c(\gamma)$ be the first cell that it does not determine and
$e(\gamma)$ the neighbourhood above it. The \emph{children} of $\gamma$ are the
nodes $\gamma\cup\{e(\gamma)\mapsto v\}$ for $v\in W$, except
\begin{enumerate}[(a)]
\item $v=\Lq$ if $c(\gamma)$ is a front cell $(t,t+1)$ with $t\le N-1$;
\item $v=\Gq$, in a class in which \cref{lem:wall} holds; in such a class the
cells $(t,1)$ are the constant $\Gq$, carry no variables and are left out of the
order $\prec$;
\item $v=\Bq$ if $\Aq$ is not a value of $\gamma$.
\end{enumerate}
The root is the empty map. A set $\mathcal U$ of nodes \emph{covers} a node
$\gamma$ if $\gamma\in\mathcal U$, or if $c(\gamma)$ exists and $\mathcal U$
covers all children of $\gamma$. We write $F\wedge\gamma$ for a formula $F$
with the unit clauses $[\delta(e)=\gamma(e)]$, $e\in\operatorname{dom}\gamma$.

\begin{lemma}[complete case split]\label{lem:split}
Let the class be invariant under exchanging the names $\Aq$ and $\Bq$, and let
$\mathcal U$ be a finite set of pairs $(\gamma,N_\gamma)$ with $3\le N_\gamma\le N$
such that the nodes $\gamma$ cover the root of the tree for $N$. If
$F^{(N_\gamma)}\wedge\gamma$ is unsatisfiable for every pair, then no rule of the
class synchronizes all lines of lengths $2,\dots,N+2$ in minimal time. If
$N\ge5$ and case \textup{(b)} does not occur in the class, then no five-state
rule of the class synchronizes all lines of lengths $2,\dots,N+1$.
\end{lemma}

\begin{proof}
Let $\delta$ be such a rule. Consider the neighbourhoods above the cells of
its half-line in the order $\prec$; if the first of them whose image is an
auxiliary state has the image $\Bq$, exchange the names $\Aq$ and $\Bq$. The new
rule is in the class and synchronizes the same lines. We show by induction that
every covered node $\gamma\subseteq\delta$ has a descendant
$\gamma'\subseteq\delta$ that occurs in $\mathcal U$. If $\gamma$ occurs in
$\mathcal U$, there is nothing to show. Otherwise $c(\gamma)$ exists; let
$v:=\delta(e(\gamma))$, which is the state of the cell $c(\gamma)$ in the
half-line of $\delta$. By (F4), $v\ne\Lq$ in case (a); by \cref{lem:wall} with
$N+2$ in place of $N$, which reaches the anti-diagonal $2N$, the cells $(t,1)$
are in state $\Gq$ and $v\ne\Gq$ in case (b); and $v\ne\Bq$ in case (c), because $e(\gamma)$ is then the
first neighbourhood of the half-line with an auxiliary image. So
$\gamma\cup\{e(\gamma)\mapsto v\}\subseteq\delta$ is a child of $\gamma$, hence
covered, and the induction applies to it; the tree is finite. For the root we
obtain a pair $(\gamma,N_\gamma)$ with $\gamma\subseteq\delta$. By
\cref{lem:Fsound}, $\delta$ satisfies $F^{(N_\gamma)}$, and it satisfies the
unit clauses of $\gamma$; this contradicts the hypothesis. For the second
statement, let $\delta$ synchronize the lines of lengths $2,\dots,N+1$; the
argument is the same with the second statement of \cref{lem:Fsound}, which
gives (F4) and $F^{(N_\gamma)}$, and without case (b).
\end{proof}

The children of a node in the tree for $N$ are among its children in the tree
for $N'\le N$, and $c(\gamma)$, $e(\gamma)$ do not depend on $N$ as long as
$c(\gamma)$ exists for the smaller value. The search program therefore works
level by level. At the level $N$ it walks the tree below its start nodes; for
a node $\gamma$ it runs CaDiCaL~3.0.1 on $F^{(N)}\wedge\gamma$ with a time limit
of $4$~s and an LRAT proof. If the answer is \textsc{unsatisfiable}, the proof
ends with the empty clause and \texttt{lrat-check} accepts it, the pair
$(\gamma,N)$ is recorded with the SHA-256 digests of the formula and of the
proof; if the solver finds a model, $\gamma$ becomes a start node of the level
$N+1$; otherwise the children of $\gamma$ are treated. A second program, which
shares no code with the first, reads the recorded pairs and checks that they
cover the root.

\begin{proof}[Proof of \cref{thm:stable}]
The second statement is the rule \texttt{stable\_generals\_2-14}. For the first,
let $\delta$ be a five-state rule with (B1) that synchronizes the lines of
lengths $2,\dots,16$. By the remark after \cref{lem:fire} we may assume that
exactly the three neighbourhoods of (B3) have the image $\Fq$. The class of the
rules with (B1) and this property is invariant under exchanging $\Aq$ and
$\Bq$, and \cref{lem:wall} holds in it.

Let $F_{14}$ be $F^{(14)}$ for this class, where (F5) consists of the unit
clauses of (B1) and of (B3) in the strong reading, of the unit clauses
$C_\infty(t,1)=\Gq$ and $C_\infty(t,i)\ne\Gq$ ($i\ge2$) for the half-line cells
with $t+i\le27$, and of three groups of redundant clauses: the statements (2b),
(2c) and (2d) of \cref{lem:fire} for the lines of lengths $n\le14$; for
$n\le14$ and $n-1\le t\le2n-4$, no two positions of the line of length $n$ at a
distance $d$ with $2\le d\le\ell+1$ are both barriers at time $t$, where $\ell$
is the largest integer with $M(\ell)<2n-3-t$; and $\PUMP(n,n')$ for
$4\le n<n'\le14$. $F_{14}$ has
$3465$ variables and $324\,753$ clauses. The rule $\delta$ satisfies it: for
(F1)--(F4) this is \cref{lem:Fsound}; \cref{lem:wall} with $N=16$ gives the unit
clauses on the half-line; the first two groups of redundant clauses are \cref{lem:fire} and the last
statement of \cref{lem:gap} for the lengths $n\le14$; and $\PUMP(n,n')$ is
\cref{lem:pumping}.

The search produced $44\,516$ nodes $\gamma$ for which $F_{14}\wedge\gamma$ is
refuted with a verified proof (\cref{tab:stable}), and one node $\gamma^\ast$,
with $11$ transitions, for which $F_{14}\wedge\gamma^\ast$ is satisfiable; its
model is the rule for $2,\dots,14$. For $\gamma^\ast$ we use the formula
$S_{15}$ with the clauses (F1) and (F2) for the lines of lengths
$2,\dots,15$, (F3) for $n\le14$, (F4) for $t\le12$, and (B1): every rule with
(B1) that synchronizes the lines of lengths $2,\dots,15$ satisfies it, by
\cref{lem:return}. CaDiCaL refutes $S_{15}\wedge\gamma^\ast$ ($3556$ variables,
$384\,704$ clauses), and \texttt{lrat-check} and \texttt{cake\_lpr} verify the
proof. The $44\,517$ nodes cover the root of the tree for $N=14$ (the checking
program walks $66\,808$ nodes). As in the proof of \cref{lem:split}, the rule
$\delta$ contains one of them. It satisfies $F_{14}$ and $S_{15}$ with the unit
clauses of that node, which contradicts the refutations.
\end{proof}

\begin{table}[t]
\centering\small
\caption{The certificates of \cref{thm:stable}. The tree was cut at depth $7$
into $341$ subtrees. ``Time'' is the CaDiCaL time on the closed nodes, and the
last column is the total size of their LRAT proofs. A node whose transitions
determine the whole half-line up to the anti-diagonal $27$ cannot be split;
four such nodes needed more than $4$~s.}
\label{tab:stable}
\begin{tabular}{@{}llrrr@{}}
\toprule
Nodes & Formula & Closed & Time & Proofs\\
\midrule
above depth $7$ & $F_{14}$ & $36$ & $<1$ s & $0.02$ GB\\
subtrees $1$ to $167$ & $F_{14}$ & $17\,196$ & $35\,384$ s & $178.1$ GB\\
subtrees $168$ to $341$ & $F_{14}$ & $27\,280$ & $57\,808$ s & $260.5$ GB\\
complete half-line, longer limit & $F_{14}$ & $4$ & $15$ s & $0.07$ GB\\
the node $\gamma^\ast$ & $S_{15}$ & $1$ & $12$ s & $0.08$ GB\\
\bottomrule
\end{tabular}
\end{table}

\begin{proof}[Proof of \cref{prop:weakimage}]
Exchanging the names $\Aq$ and $\Bq$ maps the class for $I$ onto the class for
$(\Aq\,\Bq)\,I\,(\Aq\,\Bq)$. The classes for the identity, $(\Aq\,\Bq)$,
$(\Lq\,\Gq)$ and $(\Lq\,\Gq)(\Aq\,\Bq)$ are invariant, and the exchange maps the
classes for $(\Lq\,\Aq)$, $(\Gq\,\Aq)$ and $(\Lq\,\Aq)(\Gq\,\Bq)$ onto those for
$(\Lq\,\Bq)$, $(\Gq\,\Bq)$ and $(\Lq\,\Bq)(\Gq\,\Aq)$. For each involution, (F5)
consists of the clauses of the weak reading of (B1) and of \eqref{eq:B2o}, and
case (b) of the tree does not occur. The search was run for each of the ten
involutions, from the level $3$ on (for the identity at the level $13$ only);
\cref{tab:weakimage} lists the levels at which a node was closed. Every proof
was verified by \texttt{lrat-check} and by \texttt{cake\_lpr}, and for every $I$
the nodes closed at the levels up to the last level $N'_I$ cover the root of
the tree for $N'_I$. The last levels are $13$ for the identity, $11$ for
$(\Aq\,\Bq)$, and $10$ and $11$ for $(\Lq\,\Aq)$ and $(\Lq\,\Bq)$. For an
invariant class, the second statement of \cref{lem:split} applies with
$N=N'_I$. For two classes that the exchange maps onto each other, the choice of
names in the proof of \cref{lem:split} turns a rule of one of them into a rule
of one of the two classes whose half-line is a path of the tree of that class,
and the proof of \cref{lem:split} applies in that tree; this gives the larger
of the two last levels for both. Hence $N^\ast_I\le13$ for the identity and
$N^\ast_I\le11$ for the other three. The lower bounds are
\texttt{balzer\_noborder12} and two rules found by the solver, for
$(\Aq\,\Bq)$ and for $(\Lq\,\Bq)$, which satisfy the conditions on their tables
and synchronize the lines of lengths $2,\dots,10$ (direct simulation).

For the six involutions that move $\Gq$, the formula
$\MT_5(\{2,\dots,N\};2N-2)$ with the clauses of the two conditions decides
$N^\ast_I$: it is unsatisfiable for $N=3$ if $I$ is $(\Lq\,\Gq)$ or
$(\Lq\,\Gq)(\Aq\,\Bq)$ and for $N=7$ if $I$ is one of the other four, in each
case with an LRAT proof that \texttt{lrat-check} and \texttt{cake\_lpr} verified,
and it is satisfiable for $N=2$ and for $N=6$, respectively; the rules that the
solver returns synchronize the lines of these lengths (direct simulation). The
tree search also closes these six classes, at the levels of
\cref{tab:weakimage}.
\end{proof}

\begin{table}[t]
\centering\small
\caption{The search of \cref{prop:weakimage}: one row for every level at which
a node was closed; at the lower levels the root is satisfiable. ``Sat.''
counts the nodes with a model, which are searched again at the next level.}
\label{tab:weakimage}
\begin{tabular}{@{}lrrrrrrrr@{}}
\toprule
$I$ & Level & Variables & Clauses & Closed & Split & Sat. & Time & Proofs\\
\midrule
identity & $13$ & $2608$ & $259\,620$ & $12\,296$ & $4213$ & $0$ & $20\,727$ s & $102.3$ GB\\
$(\Aq\,\Bq)$ & $10$ & $1536$ & $121\,369$ & $677$ & $258$ & $1$ & $1090$ s & $4.62$ GB\\
$(\Aq\,\Bq)$ & $11$ & $1840$ & $160\,062$ & $4$ & $1$ & $0$ & $8$ s & $0.03$ GB\\
$(\Lq\,\Aq)$ & $9$ & $1280$ & $89\,308$ & $115$ & $50$ & $12$ & $193$ s & $0.79$ GB\\
$(\Lq\,\Aq)$ & $10$ & $1536$ & $121\,369$ & $190$ & $67$ & $0$ & $290$ s & $1.25$ GB\\
$(\Lq\,\Bq)$ & $9$ & $1280$ & $89\,308$ & $32$ & $15$ & $7$ & $58$ s & $0.25$ GB\\
$(\Lq\,\Bq)$ & $10$ & $1536$ & $121\,369$ & $218$ & $81$ & $2$ & $382$ s & $1.66$ GB\\
$(\Lq\,\Bq)$ & $11$ & $1840$ & $160\,062$ & $26$ & $10$ & $0$ & $39$ s & $0.18$ GB\\
$(\Gq\,\Aq)$ & $7$ & $896$ & $42\,750$ & $1$ & $0$ & $0$ & $0.2$ s & $2.1$ MB\\
$(\Gq\,\Bq)$ & $7$ & $896$ & $42\,750$ & $1$ & $0$ & $0$ & $0.3$ s & $2.0$ MB\\
$(\Lq\,\Aq)(\Gq\,\Bq)$ & $6$ & $760$ & $27\,111$ & $1$ & $0$ & $0$ & $0.2$ s & $1.5$ MB\\
$(\Lq\,\Bq)(\Gq\,\Aq)$ & $6$ & $760$ & $27\,111$ & $1$ & $0$ & $0$ & $0.1$ s & $1.4$ MB\\
$(\Lq\,\Gq)$ & $3$ & $528$ & $3566$ & $1$ & $0$ & $0$ & $<0.1$ s & $4.8$ kB\\
$(\Lq\,\Gq)(\Aq\,\Bq)$ & $3$ & $528$ & $3590$ & $1$ & $0$ & $0$ & $<0.1$ s & $4.8$ kB\\
\bottomrule
\end{tabular}
\end{table}

Three remarks on the scope of these statements. First, Mazoyer's six-state
solution has no stable general (\cref{sec:balzer}), so \cref{thm:stable} does
not concern rules of its kind. It satisfies (B3), but with (B3) alone the same
case split is far out of reach: random descents in the tree for the lines of
lengths $2,\dots,13$ suggest some $10^{15}$ nodes, where the same estimate for
the class of \cref{thm:stable} gives $10^{6}$. Second, the proofs of
\cref{tab:stable} were checked by \texttt{lrat-check}; those of the first row
and of the last two rows, and those of a random sample of $600$ of the
$44\,476$ nodes of the subtrees, were verified by \texttt{cake\_lpr} as well,
after the formulas and the proofs had been produced again with the recorded
digests. Third, the proof of
\cref{thm:stable} says where a five-state rule with (B1) for the lines of
lengths $2,\dots,15$ would have to be looked for: such a rule has a wall on the
anti-diagonal $27$ of its half-line, $C_\infty(27-i,i)=\Gq$ for some
$3\le i\le13$. Indeed, it satisfies $S_{15}$ and every clause of $F_{14}$ except
possibly the unit clauses that use the line of length $16$, which are (F4) for
the front cell $(13,14)$ and the unit clauses $C_\infty(27-i,i)\ne\Gq$ for
$2\le i\le14$. The front states $f_t=C_\infty(t-1,t)$, $2\le t\le13$, lie in
$\{\Aq,\Bq\}$ by \cref{lem:return,lem:wall} for the lengths up to $15$, and
$f_{t+1}=\delta(f_t,\Lq,\Lq)$; so $f_{14}\in\{f_3,f_4\}$ is $\Aq$ or $\Bq$,
which settles the cell $(13,14)$. Moreover $C_\infty(25,2)\ne\Gq$, because
otherwise the line of length $15$ would have a cell in state $\Fq$ at time
$26$. If none of the cells $(27-i,i)$, $3\le i\le13$, were in state $\Gq$,
the rule would satisfy $F_{14}$, its half-line would be a path of the tree, and
the last step of the proof of \cref{thm:stable} would apply to it.
